\documentclass[11pt]{article}

\usepackage{amsmath,amssymb}
\usepackage{geometry}
\usepackage{graphicx}
\usepackage{hyperref}
\usepackage{pgfplots}
\usepackage{listings}
\usepackage{xcolor}
\usepackage{float}
\usepackage{booktabs}

\pgfplotsset{compat=1.18}
\geometry{margin=1in}

\lstdefinestyle{code}{
  basicstyle=\ttfamily\small,
  breaklines=true,
  frame=single,
  columns=fullflexible,
  keepspaces=true,
  showstringspaces=false,
  captionpos=b,
}

\title{RGB-D Object Pose Estimation and Augmented Reality Overlay in AVP World Space under Platform Constraints}
\author{}
\date{}

\begin{document}
\maketitle

% ============================================================
\section{Introduction}

This document presents the design, implementation, and evaluation of a system for estimating and visualizing the $6$D pose of a physical object in the Apple Vision Pro (AVP) world coordinate system. The system integrates external RGB-D sensing, camera calibration, ArUco-based spatial registration, and learning-based pose estimation using FoundationPose, with a visionOS user interface for configuration and visualization.

A defining challenge of this project is the lack of direct access to AVP camera feeds under the standard Apple Developer Program (ADP). At the time of development, ADEP is the only supported mechanism for accessing AVP camera data. Since ADEP access is unavailable, the system relies on external RGB-D sensors and AirPlay mirroring (UxPlay) to obtain camera observations, followed by reprojection into the AVP coordinate space. This constraint shapes the system architecture, calibration strategy, and evaluation methodology.

The document is structured as follows. Section~\ref{sec:design} describes the system design and explored alternatives. Section~\ref{sec:implementation} details the implementation of the final pipeline. Section~\ref{sec:evaluation} presents the evaluation strategy and results. The appendices provide the mathematical formulation, system diagrams, API documentation, usage instructions, AR interface description, and representative code listings.

% ============================================================
\section{Design}
\label{sec:design}

\subsection{Design Objectives}

The system is designed to satisfy the following objectives:

\begin{itemize}
    \item Estimate accurate $6$D object poses from RGB-D data
    \item Enable stable visualization of poses in AVP world space
    \item Isolate pose estimation accuracy from world alignment error
    \item Support intuitive user interaction for segmentation and calibration
    \item Remain robust under platform-imposed constraints
\end{itemize}

\subsection{Architectural Overview}

The system is decomposed into the following core modules:

\begin{itemize}
    \item \textbf{AirPlay RGB Acquisition}: UxPlay capture of mirrored AVP frames
    \item \textbf{External RGB-D Acquisition}: Intel RealSense sensor
    \item \textbf{RGB-D Alignment and Reprojection}: RealSense depth reprojected into the AVP view
    \item \textbf{Segmentation and Mask Generation}: Normalized circular ROI masks
    \item \textbf{Pose Estimation (FoundationPose)}: $6$D pose inference on RGB-D + mask
    \item \textbf{World Registration}: ArUco-based calibration with continuous tracking
    \item \textbf{visionOS User Interface}: configuration, monitoring, and overlay rendering
\end{itemize}

This modular structure allows intermediate results (e.g., camera-frame poses) to be evaluated independently of world-space visualization.

\subsection{Platform Restrictions and Their Implications}

The lack of ADEP access prevents direct use of AVP camera streams and depth sensing. Table~\ref{tab:adep_blocked} summarizes features that are inaccessible.

\begin{table}[H]
    \centering
    \caption{Features inaccessible due to lack of ADEP access.}
    \label{tab:adep_blocked}
    \begin{tabular}{p{0.45\linewidth}p{0.45\linewidth}}
        \toprule
        \textbf{Feature} & \textbf{Availability} \\
        \midrule
        AVP RGB camera feed & Blocked (ADEP) \\
        AVP depth sensing & Blocked (ADEP) \\
        Camera intrinsics & Blocked (ADEP) \\
        Direct pixel-to-world mapping & Blocked (ADEP) \\
        Native occlusion reasoning & Blocked (ADEP) \\
        \bottomrule
    \end{tabular}
\end{table}

As a result, the system must use an external RGB-D sensor and perform multi-stage reprojection with ArUco-based calibration and device-to-camera offset estimation.

\subsection{Explored Alternative Approaches}

Several alternative approaches were explored during the design phase.

\paragraph{Monocular Depth Estimation}
Transformer-based monocular depth estimation (AnyDepth v2) was evaluated as a potential replacement for external depth sensing. While visually plausible, monocular depth exhibited scale ambiguity and temporal instability when compared to RealSense depth. This approach was therefore excluded from the final system and retained only for offline experiments.

\paragraph{ROI and Occlusion Handling}
More complex segmentation strategies were explored, including freeform ROIs and active contour optimization. These methods increased interaction complexity without providing significant benefit, as FoundationPose can converge with very sparse valid pixels.

Temporal occlusion handling using Kalman filters and Gaussian processes was also evaluated. These methods introduced additional latency and tuning complexity and were excluded from the final pipeline.

\subsection{User-Centered Interaction Design}

A core design decision was to prioritize usability. Object segmentation is controlled via a simple disc-shaped ROI and a dedicated ROI overlay window rather than freehand drawing. This reduces cognitive load and enables fast iteration during calibration and evaluation.

% ============================================================
\section{Implementation}
\label{sec:implementation}

\subsection{Hardware and Software Environment}

RGB-D data is acquired using a fixed Intel RealSense sensor. The computer vision pipeline runs on a Kubuntu server using Docker containers. Development and debugging of the AR interface are performed on macOS using Xcode and the visionOS simulator, and deployments are tested on the Apple Vision Pro headset.

FoundationPose inference is executed via a Dockerized API deployed on the Kubuntu server. This work integrates the API but does not claim authorship or modification of the FoundationPose implementation.

\subsection{RGB-D Alignment and Reprojection}

RealSense depth is first aligned to the RealSense color camera frame using the vendor SDK. After alignment, each depth pixel corresponds to a ray in the color camera frame and is back-projected using the color intrinsics.

The aligned depth is then reprojected into the AVP camera frame (captured via UxPlay) using a camera-to-camera transform estimated via ArUco markers. This yields an RGB-D representation in the AVP view, which is required for the AVP-based FoundationPose pipeline.

\subsection{Segmentation and Mask Generation}

Two ROI-based masking strategies are implemented:

\begin{itemize}
    \item \textbf{AVP pipeline}: A normalized circular ROI mask defined by \texttt{/avp\_roi\_config} (center $(c_x,c_y)$ and radius $r$ in normalized coordinates). The mask is applied directly in the AVP (UxPlay) image plane.
    \item \textbf{RealSense reference pipeline}: A geometric circular ROI in the RealSense RGB view. The circle center and radius are controlled via \texttt{/rs\_roi\_config}, and the mask is visualized in the debug viewer.
\end{itemize}

The AVP ROI mask is intentionally simple and robust to illumination changes. Color-based HSV segmentation and freeform ROIs were evaluated but removed due to interaction complexity and limited benefit for FoundationPose.

\subsection{Pose Estimation}

Given aligned RGB-D data, camera intrinsics, a binary mask, and a selected PLY object model, FoundationPose produces a $6$D object pose in the corresponding camera frame. The required inputs are:

\begin{itemize}
    \item RGB image
    \item Depth map
    \item Camera intrinsics matrix $\mathbf{K}$
    \item Binary object mask
    \item Object mesh (PLY)
\end{itemize}

The output is a rigid transform $\mathbf{T}_{C \leftarrow O}$, where $C$ is either the RealSense or AVP camera frame.

\subsection{World Registration}

ArUco markers placed in the scene provide the spatial relationship between the AVP camera and the physical environment. When a marker is detected in the AVP camera frame, the pose $\mathbf{T}_{C \leftarrow M}$ is estimated via OpenCV. The AVP camera pose in world coordinates is computed using the ARKit device anchor and a calibrated camera offset:
\[
\mathbf{T}_{W \leftarrow C} = \mathbf{T}_{W \leftarrow D} \mathbf{T}_{D \leftarrow C},
\]
\[
\mathbf{T}_{W \leftarrow M} = \mathbf{T}_{W \leftarrow C}\mathbf{T}_{C \leftarrow M}.
\]

For continuous tracking, the device-to-marker transform $\mathbf{T}_{D \leftarrow M}$ is stored when the marker is visible. If the marker is temporarily occluded, the system estimates the marker pose as
\[
\mathbf{T}_{W \leftarrow M} \approx \mathbf{T}_{W \leftarrow D} \mathbf{T}_{D \leftarrow M}.
\]

This approach removes the need for a manual gizmo anchor while keeping the marker pose stable under brief occlusions. A small camera offset calibration (in meters) is exposed in the Anchor Setup window to account for device-to-camera extrinsics.

\subsection{visionOS Application}

The visionOS application provides configuration, debugging, and immersive visualization. The main window allows:

\begin{itemize}
    \item API host and port configuration
    \item Object mesh selection
    \item ROI radius visualization and ROI window
    \item Immersive space toggling
    \item Access to debug, help, and anchor setup windows
\end{itemize}

The immersive space displays world axes, the ArUco board axes, the RealSense camera pose, and the current object pose as coordinate frames and arrows.

% ============================================================
\section{Evaluation}
\label{sec:evaluation}

\subsection{Evaluation Strategy}

To isolate pose estimation accuracy from world alignment error, two pipelines are compared:

\begin{enumerate}
    \item \textbf{RS--FoundationPose (Reference)}: Pose estimated directly in the RealSense color frame using aligned RealSense RGB-D and a circular ROI mask in the RealSense view.
    \item \textbf{RS--AVP--FoundationPose (Evaluated)}: Pose estimated in the AVP camera frame using RealSense depth reprojected into the AVP view and a normalized AVP ROI mask, then transformed back to the RealSense frame.
\end{enumerate}

The RealSense-based pipeline serves as a reference due to its fixed geometry, native depth, and simple segmentation. For evaluation, the AVP-based pose $\mathbf{T}_{AVP \leftarrow O}$ is transformed into the RealSense frame using the camera-to-camera transform, yielding $\mathbf{T}^{\text{eval}}_{RS \leftarrow O}$. This is compared against $\mathbf{T}^{\text{ref}}_{RS \leftarrow O}$ from the reference pipeline. Paired requests are captured using \texttt{/foundationpose\_save\_next} and analyzed with the evaluation script in \texttt{extra\_scripts/evaluator.py}.

\subsection{Pose Error Metrics}

For each paired sample, translation error is reported as the Euclidean distance between the reference and evaluated translation vectors, and rotation error is reported as the angle of the relative rotation
\[
\Delta \mathbf{R} = \mathbf{R}_{\text{ref}} \mathbf{R}_{\text{eval}}^\top.
\]
This metric directly reflects differences in the estimated object frame independent of world-space anchoring.

\subsection{World-Space Error Interpretation}

Discrepancies observed in AVP world space are interpreted as the result of accumulated error from:

\begin{itemize}
    \item Device-to-camera offset estimation ($\mathbf{T}_{D \leftarrow C}$)
    \item Camera intrinsics and distortion estimates
    \item ArUco pose noise and marker detection artifacts
    \item Head pose drift and ARKit device anchor errors
\end{itemize}

It is therefore possible for a pose to be consistent in the AVP image plane and the RealSense frame, yet appear offset in AVP world coordinates. Debug visualization must also be used with care: rendering the debug window inside the mirrored view can create recursive marker detections and contaminate pose estimates.

% ============================================================
\appendix

\section{Mathematical Formulation}

\subsection{Coordinate Frames and Notation}

The following coordinate frames are used:

\begin{itemize}
    \item $RS_d$: RealSense depth camera frame
    \item $RS_c$: RealSense color camera frame
    \item $AVP$: AVP camera frame (AirPlay/UxPlay capture)
    \item $D$: ARKit device anchor frame
    \item $M$: ArUco marker frame
    \item $W$: AVP world frame
\end{itemize}

Homogeneous transforms are written as
\[
\mathbf{T}_{A \leftarrow B} =
\begin{bmatrix}
\mathbf{R}_{AB} & \mathbf{t}_{AB} \\
\mathbf{0}^\top & 1
\end{bmatrix},
\]
such that
\[
\mathbf{X}_A = \mathbf{R}_{AB} \mathbf{X}_B + \mathbf{t}_{AB}.
\]

\subsection{Camera Intrinsics and Distortion}

For each camera, the pinhole intrinsics matrix is
\[
\mathbf{K} =
\begin{bmatrix}
f_x & 0 & c_x \\
0 & f_y & c_y \\
0 & 0 & 1
\end{bmatrix}.
\]

Let $(X,Y,Z)$ be a point in camera coordinates, and $(x,y)$ the normalized undistorted coordinates:
\[
x = \frac{X}{Z}, \quad y = \frac{Y}{Z}, \quad r^2 = x^2 + y^2.
\]

Using a Brown--Conrady distortion model, radial distortion is
\[
x_r = x \left( 1 + k_1 r^2 + k_2 r^4 + k_3 r^6 \right),
\]
\[
y_r = y \left( 1 + k_1 r^2 + k_2 r^4 + k_3 r^6 \right),
\]
and tangential distortion is
\[
x_d = x_r + 2 p_1 x y + p_2 (r^2 + 2 x^2),
\]
\[
y_d = y_r + p_1 (r^2 + 2 y^2) + 2 p_2 x y.
\]

Pixel coordinates $(u,v)$ are then
\[
u = f_x x_d + c_x, \quad v = f_y y_d + c_y.
\]

Intrinsics $\mathbf{K}$ and distortion coefficients $\mathbf{d} = [k_1,k_2,p_1,p_2,k_3]$ are estimated via calibration using planar targets (checkerboard or ChArUco), by minimizing the reprojection error between observed and projected target corners.

\subsection{ArUco Pose and Camera-to-Camera Transform}

ArUco pose estimation yields the marker pose in the camera frame. Using OpenCV conventions,
\[
\mathbf{X}_C = \mathbf{R}_{CM} \mathbf{X}_M + \mathbf{t}_{CM},
\]
with homogeneous transform
\[
\mathbf{T}_{C \leftarrow M} =
\begin{bmatrix}
\mathbf{R}_{CM} & \mathbf{t}_{CM} \\
\mathbf{0}^\top & 1
\end{bmatrix}.
\]

For RealSense and AVP, we obtain
\[
\mathbf{T}_{RS \leftarrow M}, \qquad \mathbf{T}_{AVP \leftarrow M}.
\]

The camera-to-camera transform from RealSense to AVP is
\[
\mathbf{T}_{AVP \leftarrow RS} =
\mathbf{T}_{AVP \leftarrow M}
\left( \mathbf{T}_{RS \leftarrow M} \right)^{-1}.
\]

The inverse of $\mathbf{T}_{RS \leftarrow M}$ is
\[
\left( \mathbf{T}_{RS \leftarrow M} \right)^{-1} =
\begin{bmatrix}
\mathbf{R}_{RSM}^\top & -\mathbf{R}_{RSM}^\top \mathbf{t}_{RSM} \\
\mathbf{0}^\top & 1
\end{bmatrix},
\]
so that
\[
\mathbf{R}_{AVPRS} = \mathbf{R}_{AVPM} \mathbf{R}_{RSM}^\top,
\quad
\mathbf{t}_{AVPRS} = \mathbf{t}_{AVPM} - \mathbf{R}_{AVPM} \mathbf{R}_{RSM}^\top \mathbf{t}_{RSM}.
\]

\subsection{Depth Alignment and Reprojection}

\paragraph{RealSense Alignment}

Raw depth pixels belong to the depth frame $RS_d$ with intrinsics $\mathbf{K}_d$. After alignment, depth is expressed in the color camera frame $RS_c$.

For an aligned depth pixel $(u,v)$ with depth $z$ (meters), the 3D point in $RS_c$ is
\[
\mathbf{X}_{RS} =
\begin{bmatrix}
X_{RS} \\ Y_{RS} \\ Z_{RS}
\end{bmatrix}
=
z \, \mathbf{K}_{RS_c}^{-1}
\begin{bmatrix}
u \\ v \\ 1
\end{bmatrix}.
\]

\paragraph{Reprojection into AVP}

Given $\mathbf{T}_{AVP \leftarrow RS}$, the point in AVP coordinates is
\[
\mathbf{X}_{AVP} =
\mathbf{R}_{AVPRS} \mathbf{X}_{RS} + \mathbf{t}_{AVPRS}.
\]

Let $\mathbf{X}_{AVP} = [X_{AVP},Y_{AVP},Z_{AVP}]^\top$. Projection into AVP image coordinates is
\[
u_{AVP} = f_x^{AVP} \frac{X_{AVP}}{Z_{AVP}} + c_x^{AVP}, \quad
v_{AVP} = f_y^{AVP} \frac{Y_{AVP}}{Z_{AVP}} + c_y^{AVP},
\]
and the depth value in AVP view is $z_{AVP} = Z_{AVP}$. A z-buffer is used to resolve collisions when multiple points map to the same $(u_{AVP},v_{AVP})$.

\subsection{Device Anchor, Camera Offset, and Continuous Tracking}

The AVP device anchor provides the transform $\mathbf{T}_{W \leftarrow D}$ in the ARKit world. The physical offset between device anchor and camera is modeled as a constant transform $\mathbf{T}_{D \leftarrow C}$, yielding the camera pose
\[
\mathbf{T}_{W \leftarrow C} = \mathbf{T}_{W \leftarrow D} \mathbf{T}_{D \leftarrow C}.
\]

When the ArUco marker is visible in the AVP camera, the marker pose in world coordinates is
\[
\mathbf{T}_{W \leftarrow M} = \mathbf{T}_{W \leftarrow C} \mathbf{T}_{C \leftarrow M}.
\]

For continuous tracking, the relative transform
\[
\mathbf{T}_{D \leftarrow M} = \left( \mathbf{T}_{W \leftarrow D} \right)^{-1} \mathbf{T}_{W \leftarrow M}
\]
is stored at detection time. If the marker becomes temporarily invisible, the system estimates
\[
\mathbf{T}_{W \leftarrow M} \approx \mathbf{T}_{W \leftarrow D} \mathbf{T}_{D \leftarrow M}.
\]

% ============================================================
\section{System Architecture and Usage}

\subsection{System Pipeline Overview}

\begin{figure}[H]
\centering
\fbox{\includegraphics[width=0.95\textwidth]{placeholder_system_pipeline.pdf}}
\caption{System pipeline overview from RGB-D acquisition to AVP visualization.}
\end{figure}

\subsection{visionOS AR Interface Overview}

The AR interface is composed of several coordinated windows that support configuration, calibration, and immersive visualization. Figure~\ref{fig:ar_ui_overview} illustrates the main UI elements.

\begin{figure}[H]
\centering
\fbox{\includegraphics[width=0.9\textwidth]{placeholder_ar_ui_overview.pdf}}
\caption{Placeholder overview of the visionOS AR interface, including main configuration window, ROI window, anchor setup (camera offset calibration), help window, debug window, and immersive space.}
\label{fig:ar_ui_overview}
\end{figure}

\subsubsection{Main Configuration Window}

The main window provides access to global settings and entry points to other tools:

\begin{itemize}
    \item API host and port fields
    \item Object mesh (PLY) selection
    \item ROI overlay radius and color
    \item Toggle for immersive space
    \item Buttons to open debug, help, and anchor setup windows
\end{itemize}

\begin{figure}[H]
\centering
\fbox{\includegraphics[width=0.8\textwidth]{placeholder_main_window.pdf}}
\caption{Placeholder main configuration window showing API configuration, mesh selection, ROI controls, and navigation buttons.}
\end{figure}

\subsubsection{ROI Control and Mask Tuning}

The ROI window provides a full-screen circular overlay that helps the user place and size a disc-shaped ROI. The backend ROI mask is configured via \texttt{/avp\_roi\_config} using normalized center and radius values; the overlay serves as a visual aid for alignment and repeatability.

\begin{figure}[H]
\centering
\fbox{\includegraphics[width=0.8\textwidth]{placeholder_roi_control.pdf}}
\caption{Placeholder ROI control window showing a disc overlay for ROI placement.}
\end{figure}

\subsubsection{Anchor Setup Window}

The anchor setup window exposes three sliders for camera offset calibration (X, Y, Z in meters). Adjustments update the device-to-camera transform used for world alignment and improve the consistency of the ArUco axes overlay.

\begin{figure}[H]
\centering
\fbox{\includegraphics[width=0.8\textwidth]{placeholder_anchor_setup.pdf}}
\caption{Placeholder anchor setup window with sliders for camera offset calibration.}
\end{figure}

\subsubsection{Help and Instructions Window}

A dedicated help window summarizes the recommended calibration and usage steps:

\begin{itemize}
    \item Starting backend services
    \item Positioning the marker and RealSense camera
    \item Configuring API endpoints
    \item Adjusting ROI parameters and camera offset calibration
    \item Avoiding recursive marker detection via the debug view
\end{itemize}

\begin{figure}[H]
\centering
\fbox{\includegraphics[width=0.7\textwidth]{placeholder_help_window.pdf}}
\caption{Placeholder help and instructions window with step-by-step guidance.}
\end{figure}

\subsubsection{Debug Window}

The debug window and backend \texttt{/debug} page show synchronized multi-view outputs: AVP RGB, RealSense RGB/depth, ROI mask visualization, transformed depth, and 2D pose overlays. This view is primarily used during development and evaluation.

\begin{figure}[H]
\centering
\fbox{\includegraphics[width=0.95\textwidth]{placeholder_debug_window.pdf}}
\caption{Placeholder debug window illustrating RGB, depth, binary mask, and 2D pose overlays.}
\end{figure}

\subsubsection{Immersive Space}

The immersive space renders the AVP world, including:

\begin{itemize}
    \item World axes
    \item ArUco board axes (with continuous tracking when occluded)
    \item The estimated RealSense camera pose
    \item The current object pose as a coordinate frame and arrow overlay
\end{itemize}

Head motion is continuously streamed to the backend, and pose overlays are updated in real time.

\begin{figure}[H]
\centering
\fbox{\includegraphics[width=0.9\textwidth]{placeholder_immersive_space.pdf}}
\caption{Placeholder immersive space showing world axes, ArUco axes, and object pose overlays.}
\end{figure}

\subsection{API Endpoints}

\begin{table}[H]
\centering
\caption{Selected API endpoints used by the AVP client and debugging tools.}
\label{tab:api_endpoints}
\begin{tabular}{p{0.30\linewidth}p{0.12\linewidth}p{0.50\linewidth}}
    \toprule
    \textbf{Endpoint} & \textbf{Method} & \textbf{Description} \\
    \midrule
    \texttt{/health} & GET & Service status, calibration flags, and ROI state \\
    \texttt{/models} & GET & List available PLY models \\
    \texttt{/select\_model} & POST & Select active mesh model \\
    \texttt{/avp\_roi\_config} & GET/POST & Configure normalized AVP ROI mask \\
    \texttt{/rs\_roi\_config} & GET/POST & Configure RealSense ROI mask (pixels) \\
    \texttt{/aruco} & GET & AVP ArUco overlay and pose in camera frame \\
    \texttt{/get\_rs\_aruco\_frame} & GET & RealSense ArUco overlay and pose \\
    \texttt{/get\_avp\_mask\_frame} & GET & AVP ROI mask visualization \\
    \texttt{/get\_transformed\_depth} & GET & RealSense depth reprojected into AVP view \\
    \texttt{/get\_rs\_pose\_in\_avp} & GET & RealSense camera pose in AVP frame \\
    \texttt{/get\_foundationpose\_pose} & GET & Latest FoundationPose output \\
    \texttt{/head\_pose} & POST & Stream head/device pose from visionOS \\
    \texttt{/debug} & GET & Web debug dashboard \\
    \texttt{/mjpeg} & GET & Live MJPEG stream with selectable view \\
    \bottomrule
\end{tabular}
\end{table}

\subsection{Usage Workflow}

A typical workflow is:

\begin{enumerate}
    \item Start the FoundationPose Docker container on the Kubuntu server.
    \item Start the main CV service (UxPlay capture, RealSense, and API).
    \item Launch the visionOS application and open the main configuration window.
    \item Configure API host and port, and select the desired object mesh.
    \item Set AVP and RealSense ROI parameters via \texttt{/avp\_roi\_config} and \texttt{/rs\_roi\_config} (or the backend \texttt{/debug} page).
    \item Open the ROI window to align the visual overlay with the target object.
    \item Open the anchor setup window and adjust the camera offset so the ArUco axes align with the physical marker.
    \item Optionally open the debug window (outside the mirrored view) to verify RGB-D alignment, mask quality, and pose overlays.
    \item Enable immersive mode to visualize the object pose, world axes, and ArUco overlay in the AVP world.
\end{enumerate}

\begin{figure}[H]
\centering
\fbox{\includegraphics[width=0.95\textwidth]{placeholder_debug_endpoint.pdf}}
\caption{Placeholder visualization of the \texttt{/debug} endpoint with RGB, depth, mask, and pose overlays.}
\end{figure}

% ============================================================
\section{Code Listings}

This appendix section provides representative code excerpts for non-trivial components of the pipeline. The listings document structure and responsibilities of major modules in the current implementation.

% ------------------------------------------------------------
\subsection{RealSense RGB-D Alignment and Back-Projection}

\begin{lstlisting}[style=code,language=Python,caption={RealSense alignment and intrinsics extraction (aligned depth to color).},label={lst:rs_align_backproject}]
# realsense_client.py
profile = self.pipeline.start(self.config)

# Align depth to the color stream
self.align_to = rs.stream.color
self.align = rs.align(self.align_to)

# Build intrinsics matrix from the color stream
color_profile = profile.get_stream(rs.stream.color)
intr = color_profile.as_video_stream_profile().get_intrinsics()
self.intrinsics = np.array([
    [intr.fx, 0, intr.ppx],
    [0, intr.fy, intr.ppy],
    [0, 0, 1],
], dtype=np.float32)
\end{lstlisting}

% ------------------------------------------------------------
\subsection{ArUco Pose Estimation and Camera-to-Camera Transform}

\begin{lstlisting}[style=code,language=Python,caption={Camera-to-camera transform from simultaneous ArUco detections.},label={lst:aruco_Tavprs}]
# final_pipeline.py
with avp_pose_lock:
    T_avp_board = avp_pose_latest.get("pose_matrix")  # T_avp<-board
with rs_aruco_lock:
    T_rs_board = rs_aruco_latest.get("pose_matrix")   # T_rs<-board

if T_avp_board is not None and T_rs_board is not None:
    T_avp_rs = T_avp_board @ np.linalg.inv(T_rs_board)
\end{lstlisting}

% ------------------------------------------------------------
\subsection{Depth Reprojection into AVP View}

\begin{lstlisting}[style=code,language=Python,caption={Reproject aligned RS depth into the AVP view with distortion and z-buffering.},label={lst:depth_to_avp}]
# final_pipeline.py (simplified)
pts = cv2.undistortPoints(uv_rs.reshape(-1, 1, 2), K_rs, dist_rs)
X = pts[:, 0, 0] * z
Y = pts[:, 0, 1] * z
points_rs = np.stack([X, Y, z, np.ones_like(z)], axis=1)
points_avp = (T_avp_rs @ points_rs.T).T

img_pts, _ = cv2.projectPoints(
    points_avp[:, :3], np.zeros(3), np.zeros(3), K_avp, dist_avp
)
u_avp = img_pts[:, 0, 0].astype(np.int32)
v_avp = img_pts[:, 0, 1].astype(np.int32)

depth_avp = np.zeros((h_avp, w_avp), dtype=np.float32)
depth_avp[v_avp, u_avp] = points_avp[:, 2]  # z-buffering applied in implementation
\end{lstlisting}

% ------------------------------------------------------------
\subsection{Normalized AVP ROI Mask}

\begin{lstlisting}[style=code,language=Python,caption={Normalized circular ROI mask for the AVP view.},label={lst:avp_roi_mask}]
# final_pipeline.py
def _compute_avp_roi_mask(h: int, w: int) -> np.ndarray:
    with avp_roi_lock:
        cfg = AvpRoiConfig(**vars(avp_roi_cfg))

    if not cfg.enabled:
        return np.full((h, w), 255, dtype=np.uint8)

    cx = int(cfg.cx_n * (w - 1))
    cy = int(cfg.cy_n * (h - 1))
    r_pix = int(cfg.r_n * float(min(w, h)))

    mask = np.zeros((h, w), dtype=np.uint8)
    cv2.circle(mask, (cx, cy), max(1, r_pix), 255, -1)
    return mask
\end{lstlisting}

% ------------------------------------------------------------
\subsection{ROI Slider Mask (Reference RS Pipeline)}

\begin{lstlisting}[style=code,language=Python,caption={Slider-controlled circular ROI mask (RS reference pipeline).},label={lst:roi_mask_rs}]
# final_pipeline.py
mask = np.zeros(rgb.shape[:2], dtype=np.uint8)
cv2.circle(mask, (cx, cy), radius, 255, -1)
\end{lstlisting}

% ------------------------------------------------------------
\subsection{API Endpoints: Inference Orchestration}

\begin{lstlisting}[style=code,language=Python,caption={FoundationPose inference call with aligned RGB-D and mask.},label={lst:api_infer}]
# final_pipeline.py
rgb = cv2.cvtColor(frame, cv2.COLOR_BGR2RGB)
pose = estimate_pose(
    rgb=rgb,
    depth=depth_avp,
    mask=mask,
    K=processor.K,
    mesh_path=mesh_path,
    api_url=CONFIG["network"]["foundationpose_url"],
)
\end{lstlisting}

% ------------------------------------------------------------
\end{document}
